\chapter{Limiti operativi e troubleshooting}
\label{chap:troubleshooting}

\section{Premessa}

Ogni manuale tecnico completo dovrebbe esplicitare non solo il funzionamento ordinario
del sistema, ma anche i suoi limiti. Nel caso del progetto MAP questo aspetto e'
particolarmente importante, perche' il repository nasce in contesto didattico e unisce
moduli sviluppati con finalita' parzialmente diverse. Alcune funzioni sono molto
rifinite, altre sono presenti in forma minima, altre ancora richiedono un uso accorto
dell'ambiente di esecuzione. Questo capitolo raccoglie quindi le principali attenzioni
operative emerse dal confronto diretto tra codice e documentazione storica.

\section{Limiti funzionali da tenere presenti}

\subsection{Anteprima CSV non disponibile nella Home}

La schermata Home della GUI offre un pulsante di anteprima dataset, ma l'implementazione
corrente della preview non copre i file CSV. L'anteprima e' disponibile per
\emph{PlayTennis}, \emph{Iris} e database, mentre nel caso CSV la GUI mostra un messaggio
che segnala l'assenza del supporto. Di conseguenza, chi lavora con CSV deve verificare il
file esternamente oppure procedere direttamente all'esecuzione del clustering.

\subsection{Parser CSV volutamente semplice}

Il parser CSV non gestisce formati complessi con campi quotati contenenti virgole
interne. In aggiunta, la gestione dei valori mancanti e' limitata alla fase di inferenza
dei tipi e non corrisponde a una vera strategia di pulizia del dataset. Per evitare
errori difficili da diagnosticare, e' consigliabile pre-processare i CSV in modo che
contengano solo valori puliti, coerenti e facilmente interpretabili.

\subsection{Differenza tra GUI locale e modalita' client-server}

Una fonte frequente di confusione riguarda il ruolo del server. La GUI non necessita di
\module{qtServer} in esecuzione, mentre il client testuale si'. Se l'utente avvia
inutilmente il server prima della GUI non ottiene danni, ma introduce complessita'
superflua nell'ambiente di lavoro. Il server va considerato necessario solo per la
modalita' remota basata su \module{qtClient}.

\subsection{Copertura ridotta del client testuale}

Il client testuale non offre la stessa gamma di funzioni della GUI. Non carica CSV, non
visualizza grafici e non permette di esplorare i risultati con il livello di dettaglio
offerto dalla vista ad albero e dai dialoghi statistici. Il suo scopo e' soprattutto
quello di fornire un canale remoto minimale verso il backend.

\subsection{Visualizzazione Convex Hull}

L'inviluppo convesso viene disegnato solo per cluster con almeno tre punti e viene reso
come bordo, non come area piena. L'assenza del bordo in cluster molto piccoli non e'
quindi un malfunzionamento, ma una conseguenza diretta dell'algoritmo geometrico usato
per il calcolo.

\section{Problemi di compilazione ed esecuzione}

\subsection{Maven o JavaFX non disponibili}

Se \module{make run-gui} o \module{make gui} falliscono, la prima verifica da fare
riguarda la disponibilita' di Maven e della toolchain Java corretta. Il modulo
\module{qtGUI} richiede un ambiente coerente con il \module{pom.xml}; usare una versione
di Java troppo vecchia puo' produrre errori di compilazione o di risoluzione dei moduli.

\subsection{Driver JDBC mancante}

I moduli compilati con Makefile si aspettano di trovare il driver in
\filepath{qtServer/JDBC/mysql-connector-java-9.5.0.jar}. Se questo file viene spostato o
rimosso, il server e i test che dipendono dal database possono fallire con errori
relativi al caricamento del driver MySQL. In questi casi occorre ripristinare il file
nella posizione attesa o aggiornare il classpath nei Makefile.

\subsection{Errore di porta occupata}

Quando il server viene avviato su una porta gia' in uso, la classe \module{MultiServer}
termina con un errore di binding. La soluzione piu' semplice consiste nel selezionare una
porta alternativa:

\begin{commandblock}
  make run-server PORT=9090
  make run-client IP=localhost PORT=9090
\end{commandblock}

L'importante e' mantenere allineati i parametri di client e server.

\section{Problemi di connessione al database}

\subsection{Credenziali errate o database assente}

La maggior parte dei problemi JDBC deriva da configurazioni incoerenti: database non
creato, utente non autorizzato, password errata o tabella inesistente. La GUI offre un
pulsante di test connessione proprio per intercettare queste situazioni prima dell'avvio
del clustering. In modalita' client-server, invece, l'errore emerge solo al momento del
caricamento della tabella.

\subsection{Schema non adatto al clustering}

Anche quando la connessione riesce, il contenuto della tabella puo' non essere adeguato.
Colonne solo tecniche, dati sporchi o tipi SQL inattesi possono produrre dataset poco
significativi o errori in fase di caricamento. Poiche' il backend ignora le colonne
auto-increment e riduce i tipi a \emph{string} e \emph{number}, e' opportuno progettare
le tabelle come vere tabelle di feature, non come tabelle operative piene di metadati.

\section{Problemi su file ed esportazioni}

\subsection{File \module{.dmp} non apribile}

Se un file \module{.dmp} non viene aperto correttamente, il primo controllo riguarda il
nome fornito al client o il percorso selezionato nella GUI. Nel client testuale va
inserito il nome senza estensione; nella GUI, invece, il file chooser gestisce
direttamente l'estensione. Un secondo caso possibile riguarda i file legacy, che possono
essere caricati ma non sempre consentono la ricostruzione completa del dataset.

\subsection{Esportazione ZIP o report fallita}

L'export ZIP dipende dalla creazione di file temporanei e dalla possibilita' di scrivere
nella directory di destinazione. Se il sistema segnala errori I/O, la causa piu' comune
e' un problema di permessi sul percorso scelto oppure la presenza di un file gia' aperto
da un altro programma. In genere conviene riprovare su una directory semplice e locale.

\subsection{PNG con assi poco leggibili}

Quando si selezionano attributi discreti per gli assi, la visualizzazione usa codifiche
ordinali numeriche. Il grafico resta corretto dal punto di vista tecnico, ma richiede una
interpretazione qualitativa: le coordinate 0, 1, 2 non sono misure continue, bensi'
etichette numeriche associate ai valori simbolici osservati. Per grafici piu' intuitivi
e' preferibile usare attributi continui, come quelli del dataset \emph{Iris}.

\section{Buone pratiche operative}

Per un uso affidabile del sistema e' consigliabile adottare alcune regole semplici. Prima
di tutto, scegliere dati puliti e coerenti con il modello di clustering. In secondo
luogo, separare mentalmente i due scenari d'uso: GUI locale da un lato, client-server
dall'altro. Infine, salvare i risultati importanti in formato \module{.dmp} oltre che
esportarli in CSV o TXT, in modo da poterli riaprire in futuro senza dover ripetere il clustering.

Queste indicazioni non eliminano tutti i possibili problemi, ma riducono sensibilmente la
probabilita' di errori evitabili e aiutano a interpretare correttamente i messaggi
prodotti dall'applicazione.
